% page 551 du fac-simile (image p-550 du scan)
% bloc 170.2 x 242.0 mm ; marge gauche 31.8 mm
\pgc{10.160mm}{0.000mm}{
\l{\cel{53}543}
\l{}
\l{\sou{sordino}. I. (\sou{muziko}). (à) Espineto, di qua la kordin vibrigas}
\l{\cel{3}mikra peci de ligno envelopita en drapo, por moderar la soni ;}
\l{\cel{3}(b) Kono de kartono, kun truo ye la bazo, quan onu lokizas}
\l{\cel{3}en la funelo di chaskorno - mikra tubo ek ligno quan onu loki-}
\l{\cel{3}zas en la parto infra di la trumpeto, por moderar la soni ;}
\l{\cel{3}(c) Mikra peco de ligno, de ivoro, e c., di qua la formo esas}
\l{\cel{3}olta di pektilo kun tri denti kavigita, quan onu inkastras sur}
\l{\cel{3}la tresto di la violini, alti, violonceli, kontrabasi, por}
\l{\cel{3}diminutar la vibri di la framo e moderar la soni. - II. (\sou{metaf.})}
\l{\cel{3}\sou{Agar}\cel{1}sen bruiso por ke nulu remarkez ta ago. - DEFIS.}
\l{}
\l{\sou{sorito.(logiko}). Serio de propozicioni, interligita tale ke la}
\l{\cel{3}atributo di la unesma divenas la subjekto di la duesma, e tale}
\l{\cel{3}pluse - la lasta esas implicite inkluzita en la unesma se la}
\l{\cel{3}rezono esas korekta. - DEFIRS.}
\l{}
\l{\sou{sor}to. Eso-maniero di persono, di kozo (kompare a to di quo la}
\l{\cel{3}naturo esas analoga). DEFIRS.}
\l{}
\l{\sou{sospirar}. (\sou{netrans.}) Agar ta expiro profunda e pasable duroza}
\l{\cel{3}quan onu lasas eskapar kande onu opresesas da sento o da sen-}
\l{\cel{3}timento chagreniganta. - FIS.}
\l{}
\l{\sou{sovaja}. I. Qua vivas exter la socii civilizita. - II. (\sou{Pri planto})}
\l{\cel{3}Qua kreskas sen kultiveso. - III. (\sou{pri loko)}. Adube la homo}
\l{\cel{3}ne venas, ne efikas por agar o facar. - IV. (\sou{metaf.)}\cel{1}Qua ne}
\l{\cel{3}volas relatar kun sua kunhomi. - EFIS.}
\l{}
\l{\sou{sozio.}\cel{1}Persono quan onu konfundas ad altra, tante grandan esas lia}
\l{\cel{3}intersimileso. - DFI.}
\l{}
\l{\sou{spaco.}\cel{1}I. Extensajo de ca a ta punto, vakua ye korpi solida.}
\l{\cel{3}- II. (1) Extensajo ideala, egardata quale se lu kontenus omna}
\l{\cel{3}extensaji reala, omna korpi qui existas e quin la spirito kon-}
\l{\cel{3}ceptas kom posibla. (2) Parto di ta extensajo ideala. - EFIS.}
\l{}
\l{spado. Garden-implemento, ek fero-plako larja e tranchiva, muntita}
\l{\cel{3}an la extremajo di mancho longa, quan onu uzas por fendar tero}
\l{\cel{3}e kulbutigar la glebi. - DE.}
\l{}
\l{\cel{1}\sou{spadico}. (\sou{bot.)}\cel{2}Infloresenco forme di kastono, od asembluro de}
\l{\cel{4}folii sen-pedikula sur axo komuna, qua havas flori envelopita}
\l{\cel{4}da granda brakteo (spato). -}
\l{}
\l{\cel{1}\sou{spano}. Rekorturo plu o min dina quan onu deprenas (1) per raboto,}
\l{\cel{4}cizelo, kande onu fasonas ligno-peco ; (2) per utensilo tran-}
\l{\cel{4}chiva kande onu fasonas peco de stalo, de fero, de kupro, sur}
\l{\cel{4}torno-mashino, frezo-mashino, boro-mashino. - D.}
\l{}
\l{\cel{1}\sou{spanielo}. (\sou{zool.}) Chaso-hundo de fonto Hispana, di qua la pili}
\l{\cel{4}esas longa, e la oreli pendanta. - DEF.}
\l{}
\l{\cel{1}\sou{spa}nto. (navig.)Parto di la korpo di navo, varango qua formacas}
\l{\cel{4}un ek la kosti di la framo. - DR.}
}
